\begin{textsample}{NEG Dim 3 – vem\_ed\_06 – Score 5.00 – vem\_ed\_06/t023.txt}  \label{ex:f3_neg_060}
Caminhos com Viseu Realizar intercâmbios virtuais nas Fatecs, empregando a língua portuguesa.
Com essa perspectiva em mente, Osvaldo Succi Jr., coordenador dos Projetos Colaborativos Internacionais (PCIs), fez uma série de conversas ao longo de 2020 com representantes do Instituto Politécnico de Viseu (IPV), em Portugal, para iniciar PCIs com as Fatecs.
A instituição pública portuguesa oferece cursos técnicos de graduação e pós-graduação, em diversas áreas, como saúde, educação, agrária, tecnologia e gestão.
Em 22 de fevereiro de 2021, foi realizada uma reunião na plataforma Teams, para delinear caminhos possíveis para as colaborações luso-brasileiras.
Do IPV, participaram Susana Amante (professora de Língua Inglesa), António Figueiredo (atuante na área de Gestão, no Departamento de Madeiras), Joaquim Antunes (com intervenção em áreas de Marketing e Turismo), Pedro Reis (Finanças), Jorge Manuel Martins (Diretor do Departamento de Madeiras) e João Luís Pereira.
Do lado brasileiro, esteve presente a equipe dos PCIs, além de Ricardo Sérgio Neiva Nóbrega (coordenador do curso de Comércio Exterior da Fatec Indaiatuba) e Carlos Augusto Amaral Moreira (professor da Fatec Americana).
Um primeiro passo proposto pelos professores é a criação de projetos-piloto para, por exemplo, aferir diferenças culturais em ambientes de negócio no Brasil e em Portugal.
Após esse encontro virtual, Susana Amante, interlocutora da aproximação entre as instituições, concedeu o seguinte depoimento: Os projetos colaborativos internacionais constituem uma mais-valia para uma aprendizagem ativa, num processo construtivo em que os vários intervenientes – neste caso do IPV, em Portugal, e de Fatecs do Centro Paula Souza, no Brasil – se articulam para que, no âmbito das suas Unidades Curriculares, se envolvam em aprendizagens em que refletem acerca das suas próprias perspectivas e experiências, bem como sobre a forma como colegas, com outras culturas e formas de pensar, percepcionam o mundo que os rodeia e como compreendem e aplicam os conteúdos trabalhados.
Muito mais do que uma ênfase no individual, o foco é no coletivo.
A aprendizagem colaborativa implica um olhar sobre os processos para a resolução de desafios em que cada elemento contribui para uma equipa interdisciplinar e, muitas vezes, transdisciplinar, com o objetivo comum de resolução de problemas cada vez mais globais, porque somos todos cidadãos do mundo.
Neste sentido, o IPV congratula-se com esta parceria e espera que estes primeiros projetos-piloto venham a dar lugar a muitos outros PCIs produtivos, em anos vindouros.

% matched lemmas: 
\end{textsample}
